\section[第2章\q 想买却买不起的烦恼]{第2章\\ 想买却买不起的烦恼}
\label{sec:ch2}

一走出33楼的电梯，柔伊立即被密密麻麻的工作行程塞满了，和平常每个周一早上没两样。春季主题的截稿期是下周五，办公室的每位同事都卯足了全力在拚进度。各式文章、各式作者背景、还有照片标题，争相引起柔伊的关注：在厄瓜多尔尔山脉骑自行车、在巴尔干品葡萄酒、知名旅人图文并茂的文本……她的工作就是把这许多内容变成使人眼睛一亮的杂志期刊。

柔伊就在位于世贸中心的一家杂志社上班。他们也管这栋大楼叫自由之塔。这常常让柔伊觉得讽刺，因为尽管她喜欢这份忙碌的工作，但在办公大楼围墙内消磨的时光，很难被称为自由。对能在这个职位上做事，她心存感恩，但是为此付出的时间，很难与读者想像的编辑薪水成比例。

还有更讽刺的是，身为旅游期刊的副主编，27岁的她从未踏出美国国门。甚至没有走到密西西比河以西，她连一本护照都还没申请。

一个没有出过远门的旅游编辑。

她将笔电放到桌上打开，登录杂志社的内部网页后，手指在键盘上飞舞，进入工作状态。

柔伊陷入一阵忙碌漩涡：标错的截稿日、临要上稿前内容改变、把文笔平凡的稿子整理得人模人样。她一碰到键盘的她就是这样，她会跟上这个空间的节奏，把之前的犹疑不定全抛开。

\subsection{收支只能打平，有时还会失衡}

“我们还没人肚子饿啊？”

柔伊从椅子上直起身，转转脖子，放松一下身体。什么？已经过下午1：00了？她看到老板正越过半高的办公室隔板上看她。

“虽然是纸上旅游的编辑，也是要吃东西的吧？”老板又加了一句。

曼哈顿下城区处在社会地位的上流，老板芭芭拉的衣着让她好像是个小镇游客，与杂志社内其他衣着时髦的同事不大一样。和柔伊也是截然不同。但是芭芭拉极其聪颖，了解事物底下的内核。柔伊猜想，正是这个特质让芭芭拉成为了不起的编辑部总监。

柔伊6年前开始在这里工作，当时就是芭芭拉负责面试，她们一见面就很投缘。芭芭拉对手下编辑期望很高，要求也很高，却严而不厉。你不是因为怕她而把工作做好，把工作做好是因为不想让她失望。

柔伊从来没让芭芭拉失望过。她是个杰出的编辑，工作能力非常好。

“饿昏了。”柔伊说。把笔电切到休眠状态，跟着芭芭拉进电梯到搂上餐厅。

公司餐厅可以俯瞰曼哈顿区和哈德逊河，自由女神像也尽收眼底。极简风的装修搭配开放式空间，员工自助餐厅看似曼哈顿区的高档餐厅。柔伊刚来工作时，时常在餐厅巧遇名流，她还要装出一副不以为意的样子。

芭芭拉自备一个涂有亮漆的午餐盒，她小心翼翼地打开来时，柔伊正排到了午餐队伍的后方。她选了鸡肉沙拉搭配藜麦、盐烤西班牙杏仁，还有叶菜。在吃沙拉时，她和芭芭拉说起了她正在编辑的文章，但她不喜欢道人是非，短短两句后，她就不再多说了。

接下来是短暂的沉默。芭芭拉吃着三明治，看了看柔伊。过了一会儿，芭芭拉终于开口：“嗯……你看起来……有点失神。一切都还好吗？”

这就是芭芭拉善体人意的地方。柔伊正在竭力想忘记早上莫名的情绪，她的老板却感受到了。柔伊深深吸了一口气，然后吐了出来。她不知道该从哪里说起，因为她自己也还没弄明白。

“说出来，你可能会觉得奇怪。”柔伊开口说道。

芭芭拉咬了一口三明治，点点头，仿佛在说，接着说吧！

“我每天早上在搭地铁的路上会去一家咖啡厅，就在威廉斯堡文艺区那里。”柔伊描述咖啡厅所在地时，芭芭拉又点了点头。

“海伦娜咖啡（Helena’s Coffee）。”

“你知道那家啊？”

芭芭拉从三明治上方看了柔伊一眼后说：“然后？”

“好吧！柔伊说。“咖啡厅后墙上挂着裱框照片。我是说，有好多照片，整片墙都是，其中有一张特别不同寻常。”

可以想见，柔伊在点拿铁和马芬蛋糕时，排了长长的队伍。海伦娜咖啡就是那样一处地方，点心是现做的，咖啡特别可口，而张贴在墙上的照片，也总叫人过目不忘。

柔伊描述了一下那张照片后，又接着去吃沙拉。

“然后？”芭芭拉等了一下，问道。

“然后，我就不知道啊！我总在想着那张照片。我也不确定什么原因。”

柔伊是下标题高手，语句通常直接了当，但是今天她竟然说出不知道、不确定之类的话。

“你想要买下来。”

柔伊叹了一口气。那确实是她想要的。

画面非常简洁：晨光下的海边小村庄，早上第一道阳光投下琥珀色的光影，前方是准备出海的渔船。这个时光他们称之为黄金时刻，太阳刚刚升起，阳光绯红几近透明。对柔伊来说，这个片刻非常神奇，有些无形的能量正在蓄势待发。

照片尺寸非常大，大概有121公分宽、90公分高。尽管如此，有些细节她还是没有看清楚，因为时间太短，她不及细细研究。每天一早，她就匆匆离开公寓（说来其实都稍晚了一点），冲去咖啡厅点一杯拿铁和马芬蛋糕，之后快步走到车站，等往曼哈顿区的L线地铁。去了店里取早餐时她没空打量周围。

可是尽管总是匆匆忙忙，墙上的照片有时还是会引起她的注意。今天早上，她就多停留了半分钟的时间，走近了两步。所以，细节没看清楚。

当时她只是觉得，如果可以，她也想在起居室的墙壁上挂一张同样的照片。尽管所谓的起居室不过是一个长条的空间，那里还兼做她的餐厅和居家工作室。柔伊和一位室友同住一间很小的公寓，里面没什么值得一看的“风景”。但是，这张海滨的早晨却可以让陋室生辉。

“也未必是我想拥有它。只是……”只是什么？照片挑起了柔伊的某些思绪，到底是什么她一下说不上来。“我也解释不清楚。”她摇了摇头，好像要把理不清的念头甩到一边去似的，“我甚至不知道那张照片是否在出售。再说，即便是……”

芭芭拉和柔伊齐声说：“我也买不起。”

如果柔伊的生活是一首歌，这五个字就是不断重复的副歌，主歌可能很激励人心、很有冒险精神，很有想法，比方说，我想回到校园，我想去美国西南部走一走，我想去欧洲一趟，我想有一间真正的卧室，可以做瑜伽、可以写点东西。但是歌曲总在同样的副歌中结束：但是，我买不起。

问题在于，她确实买不起。布鲁克林区不比曼哈顿，但是依然不便宜。外加她还有学生贷款，那好像一袋沉甸甸的砖头，压在她身上。优点是她生活在纽约，她不需要以车代步，即便有只怕现在也被收回了。轿车？照这样下去，只怕到了夏天连她的自行车都要被收回。

柔伊审美观不错，文本功底也好，但是说到数字，那就不是她的长处了。在财务管理方面，她向来有点问题。她妈妈曾提议先列预算来节省开支。可是预算这二字对柔伊来说，简直是字典上最困难的。显然，这招没用。

在工作上，柔伊很有组织，很有效率，但对于自己的财务，她就变成一个无法自律的人。所以就成了现在的样子。已经3月了，她还在分期支付去年圣诞节刷卡买给众亲友的礼物。如果她花点时间整理收据，说不定还有几张是去年圣诞节的帐单。在柔伊的帐本上，永远是应付帐款、应付帐款、应付帐款。

说来，柔伊相当喜欢她的工作，表现也非常好，但同时也必须承认，她所赚的与花费只能刚好持平。事实上，这个平衡也常常被打破，出现捉襟见肘的窘境。柔伊有时觉得自己可以代言一台付帐的活机器。

假如那张照片是要出售的，不管标价是500美元、800美元，还是1,000美元，这样大笔的现金都是她没有的，也是她无法一时兴起就花掉的。

芭芭拉的声音打破了她的沉默。“你应该和亨利谈谈。”

“亨利？”

“就是你在咖啡厅遇见的老先生啊！他通常早上会在店里煮咖啡。他就是亨利。”

柔伊回想了一下，才想起来芭芭拉在说哪位。“你是说在海伦娜？你认识那里的咖啡师？”

芭芭拉站起身来，阖上午餐餐盒。“认识他好几年了。你该走进去，和他谈一谈。他看事情……”芭芭拉停顿了一下，“他看事情的视角很不一样。”

“和咖啡师谈一谈？”柔伊问。“要谈什么？”

芭芭拉露出了她的招牌表情，一切了然于心的样子。“就是找他谈谈，告诉他你喜欢这幅画。看看他怎么说。”

柔伊皱起眉头。

“相信我。”芭芭拉说。“他是个宝山。”

“他就会帮我吗？真的？比如买一张乐透？”

芭芭拉耸肩。“也未必是买乐透。但是你自己常说的，我买不起。那件事困扰着你，我没说错吧？”

柔伊不再说话了。她没说错，她可是芭芭拉！

“那就对了，”芭芭拉接着说。“就做点什么吧，去和亨利谈谈。”

\subsection{牺牲周末换来好薪水，自己开心吗？}

回到办公室位置上之后，柔伊心里觉得有些愧疚。她没有告诉芭芭拉真正困扰她的是什么，不是照片。还有其他。

一个去经纪公司的工作机会。

上上周五，她和大学室友洁西卡碰面。两杯酒下肚之后，洁西卡告诉她，有一家媒体经纪公司正在征人。“你那么拚，柔伊，”洁西卡说：“又聪明，文笔又好，又有人缘，你是最佳人选。”

因此上周，柔伊就找了一天，过去经纪公司面试。当天晚上，洁西卡就打电话给她说，根据她听来的消息，柔伊是最看好的。“你知道有多少人去面试吗？你把别人都比下去了。”而且，洁西卡的话很快就得到了证实，到周五晚上，她就收到经纪公司的电话，确认她是第一人选。看来只要柔伊愿意，新工作就是她的，薪水要比现在好上许多。她自然也知道，以后的责任更重，压力更大，这不是柔伊想要的。但好薪水可以解决她现在的难题。

周末时，她和妈妈谈到这件事，妈妈却不以为然。“喔，小柔，”妈妈说：“要为拥有的而开心。再说了，高薪就可以让你更快乐吗？”

金钱不能带来快乐，从小到大，柔伊听过这句话多少遍了？

爸爸把电话接了过去，这还真不寻常。“想一想吧，柔伊。”老爸说。柔伊知道这话的意思——我不会说你现在就该立即换工作，但或许你也是可以胜任这份新工作。

柔伊的父亲之前是承包商，收入丰厚。后来，因为健康出了问题，只能去一家建材公司做内勤人员。薪水当然不能与之相比，柔伊怀疑乐趣也相对减少很多，但他们的日子也照样过得下去，只是母亲近来更疲累。为拥有的而开心？不能说她的父母不开心，但是，他们可算得上真正的开心吗？

还有，柔伊自己开心吗？

她又想起早上在纽约之眼附近看到的那张照片，一条停泊在沙漠中的船。如果你不知道你要去那里……

经纪公司给柔伊一周的时间考虑，做最后决定，和办理离职手续。也就是说，如果柔伊想要这份工作，这周五就应该答复他们了。之后，她和洁西卡每周五的周末喝一杯时间，就可以好好祝贺一下。

另一个选择就是留在原来的公司，等着升迁。同时，再兼一些写作或编辑的外稿，在本就要加班的工作之外，再加一点班，晚上周末也不得闲。后者，自然也不是柔伊想要的。

但是，难道她还有其他选择吗？
